\documentclass[a4paper,10.5pt]{article}
\usepackage[utf8]{inputenc}
\usepackage[T1]{fontenc}
\usepackage[french]{babel}
\usepackage[left=1cm,right=1cm,top=.5cm,bottom=0cm]{geometry}
\usepackage{enumitem}
\usepackage{hyperref}
\usepackage{xcolor}
\usepackage{titlesec}
\usepackage{fontawesome5}
\usepackage{lmodern}
\usepackage[normalem]{ulem}

% --- Couleurs Airbus ---
\definecolor{primary}{RGB}{0, 51, 102}
\definecolor{secondary}{RGB}{80, 80, 80}

% --- Configuration des liens ---
\hypersetup{
    colorlinks=true,
    linkcolor=primary,
    filecolor=primary,
    urlcolor=primary,
}

% --- Formatage des titres ---
\titleformat{\section}{\large\bfseries\color{primary}\uppercase}{}{0em}{}[\titlerule]
\titlespacing{\section}{0pt}{8pt}{5pt}

% --- Commandes personnalisées ---
\newcommand{\resumeItem}[1]{\item\small{#1}}
\newcommand{\resumeSubheading}[4]{
  \vspace{-1pt}\item
    \begin{tabular*}{0.97\textwidth}[t]{l@{\extracolsep{\fill}}r}
      \textbf{#1} & \textit{\small #2} \\
      \textit{\small #3} & \textit{\small #4} \\
    \end{tabular*}\vspace{-5pt}
}

\begin{document}

% --- EN-TÊTE ---
\begin{center}
    {\Large \textbf{Loïc Aron MBASSI EWOLO}} \\ \vspace{4pt}
    \textbf{\large Stage : Environnement d'aide à l'ingénierie basé sur l'IA} \\
    \textit{\small Élève Ingénieur en 2ème année, double diplôme ENSEA} \\ \vspace{4pt}
    \small
    \faMapMarker\ Cergy, France (Mobile) \quad
    \faPhone\ (+33) 7 59 52 33 98 \quad
    \faEnvelope\ \href{mailto:loicaron.mbassiewolo@ensea.fr}{loicaron.mbassiewolo@ensea.fr} \\ \vspace{2pt}
    \faLinkedin\ \href{https://www.linkedin.com/in/mbassi-loic-3942a9246/}{@mbassi-loic} \quad
    \faGithub\ \href{https://github.com/Nameless0l}{@Nameless0l} \quad
    \faGlobe\ \href{https://mbassiloic.tech}{mbassiloic.tech}
\end{center}

\vspace{-8pt}

% --- PROFIL ---
\noindent\small{Élève-ingénieur en double diplôme à l'\textbf{ENSEA}, spécialisé en \textbf{IA} et \textbf{Génie Logiciel}. 
Expérience concrète en orchestration d'\textbf{agents IA} (Google ADK, RAG, LangChain) et en conception d'outils d'aide à l'ingénierie (génération de code depuis UML). 
Lauréat de 3 hackathons nationaux. Motivé à contribuer à l'intégration de l'IA dans les processus d'\textbf{ingénierie système} chez Airbus.}

\vspace{4pt}

% --- FORMATION ---
\section{Formation}
\begin{itemize}[leftmargin=*, label={.}, nosep]
    \item \textbf{Double diplôme : École Nationale Supérieure de l'Électronique et de ses Applications (ENSEA)} \hfill \textit{2025 -- Présent} \\
    \small{2ème année - Spécialité Intelligence Artificielle et Traitement du Signal}
    \vspace{2pt}
    \item \textbf{École Nationale Supérieure Polytechnique de Yaoundé — Génie Informatique} \hfill \textit{2021 -- Présent} \\
    \small{Niveaux 4 \& 5 : Architectures logicielles, Design Patterns, Systèmes distribués, IA/ML} \\
    \small{Niveaux 1 à 3 : Algorithmique, programmation C, bases de données, UML}
\end{itemize}

% --- COMPÉTENCES TECHNIQUES ---
\section{Compétences Techniques}
\begin{itemize}[leftmargin=*, nosep]
    \item \small{\textbf{Intelligence Artificielle :} Python, PyTorch, TensorFlow, \textbf{LangChain}, Google ADK, \textbf{RAG}, Transformers, NLP}
    \item \small{\textbf{Développement :} \textbf{C\#/.NET} (notions), PHP/Laravel, Java, TypeScript, \textbf{Design Patterns}, REST APIs}
    \item \small{\textbf{Outils \& DevOps :} Git/GitHub, \textbf{Docker}, CI/CD, tests unitaires, méthodologie \textbf{Agile/Scrum}}
    \item \small{\textbf{Ingénierie :} UML, Architecture logicielle, Documentation technique, Microservices}
\end{itemize}

% --- PROJETS IA & INGÉNIERIE ---
\section{Projets — IA \& Outils d'aide à l'ingénierie}
\begin{itemize}[leftmargin=*, label={}]

    \item \textbf{Système Multi-Agents Agricole — Orchestration IA} \hfill \textit{2024 -- 2025, recherche}
    \vspace{-2pt}
    \begin{itemize}[leftmargin=0.4cm, nosep, label=\tiny$\bullet$]
        \item \small{Conception et orchestration de \textbf{5 agents IA} spécialisés avec mécanismes de résolution de conflits}
        \item \small{Implémentation d'un système \textbf{RAG} connecté à une base de connaissances structurée}
        \item \small{Architecture modulaire facilitant l'ajout de nouveaux \textbf{cas d'usage} — Google ADK, \textbf{LangChain}, Docker}
    \end{itemize}
    \vspace{4pt}

    \item \textbf{Laravel API Generator — Génération de Code Assistée par LLM} \hfill \textit{2024 -- Présent, Open Source}
    \vspace{-2pt}
    \begin{itemize}[leftmargin=0.4cm, nosep, label=\tiny$\bullet$]
        \item \small{Outil générant une architecture backend complète depuis des diagrammes \textbf{UML/XML}}
        \item \small{Intégration de \textbf{LLM} pour analyser les spécifications et proposer des optimisations}
        \item \small{Parsing avancé et métaprogrammation pour automatiser les tâches d'\textbf{ingénierie logicielle} — \href{https://github.com/Nameless0l/laravel-api-generator}{\faGithub}}
    \end{itemize}
    \vspace{4pt}

    \item \textbf{UROP 2024 — Analyse ECG par Deep Learning (Collaboration Google DeepMind)} \hfill \textit{2024}
    \vspace{-2pt}
    \begin{itemize}[leftmargin=0.4cm, nosep, label=\tiny$\bullet$]
        \item \small{Modèle de détection automatique d'anomalies cardiaques — TensorFlow/Keras, Computer Vision}
        \item \small{Extraction de features et prétraitement de données médicales complexes}
    \end{itemize}

\end{itemize}

% --- EXPÉRIENCES PROFESSIONNELLES ---
\section{Expériences Professionnelles \& Recherche}
\begin{itemize}[leftmargin=*, label={}]

    \resumeSubheading
      {Stage de Recherche à distance — MILA}{Juin -- Août 2025}
      {Institut Québécois d'Intelligence Artificielle}{Québec, Canada}
      \resumeItem{Grokking (OpenAI) : Généralisation des réseaux de neurones après sur-apprentissage}
      \begin{itemize}[leftmargin=0.4cm, nosep, label=\tiny$\bullet$]
        \item \small{Reproduction d'expériences SOTA, implémentation de pipelines de tests sous \textbf{PyTorch}}
        \item \small{Analyse mathématique de la convergence des modèles et documentation des résultats}
      \end{itemize}
    \vspace{2pt}

    \resumeSubheading
      {Développeur Backend (Fintech) — CSC}{Oct. 2024 -- Jan. 2025}
      {Cameroon Software Company}{Yaoundé, Cameroun}
      \resumeItem{Conception d'une plateforme de transactions financières sécurisée}
      \begin{itemize}[leftmargin=0.4cm, nosep, label=\tiny$\bullet$]
        \item \small{Architecture backend scalable (Laravel, MySQL, \textbf{Docker}), \textbf{CI/CD}, tests unitaires}
        \item \small{Travail collaboratif en équipe \textbf{Agile/Scrum}, documentation technique et revues de code}
      \end{itemize}

\end{itemize}

% --- HACKATHONS & LEADERSHIP ---
\section{Hackathons \& Leadership}
\begin{itemize}[leftmargin=*, nosep]
    \item \small{\textbf{Hackathons :} 1\textsuperscript{re} place IndabaX Africa 2025 (IA Santé) | 1\textsuperscript{re} place CITS National \& Mountain Hub Regional (2024)}
    \item \small{\textbf{Leadership :} Lead AI Cell ENSPY (+50\% membres), Animation workshops ML/LLM, Articles techniques Medium}
\end{itemize}

% --- CERTIFICATIONS ---
\section{Certifications \& Langues}
\begin{itemize}[leftmargin=*, nosep]
    \item \small{Supervised Machine Learning (DeepLearning.AI / Andrew Ng) | Python for Data Science (IBM)}
    \item \small{\textbf{Langues :} Français (Langue maternelle), \textbf{Anglais} (Professionnel, documentation technique)}
\end{itemize}

\end{document}